\bigskip\Needspace{12\baselineskip}
\pdfbookmark[1]{Gauss to Bode. 3 March 1803}{bode18030303}
\begin{center}
GAUSS TO BODE.\par
\medskip\hrule\medskip
Astronomisches Jahrbuch für 1806. pp.\ 179—180. Berlin 1803.\par
\medskip\hrule
\end{center}\medskip
\begin{flushright}Braunschweig, 3 March 1803.\end{flushright}

Your — will probably already have received the important news from Mr.\ \textsc{Harding} himself that \textsc{Olbers}’s planet was found again so early at Lilienthal. You probably already have Dr.\ \textsc{Olbers}’s observations as well. I hope you will welcome a fuller account of the results as well. The weather here has been very bad since 23 Feb., and now there is the moonlight as well. Observations will therefore probably remain interrupted a little longer.

Here is the comparison of \textsc{Olbers}’s observation with my elements V, from which the ephemeris in the December issue of the M.\ C. was calculated.
\begin{center}
\begingroup\small\setlength{\tabcolsep}{3pt}
\begin{tabular}{l|r|r|r|r|r}
& &\multicolumn{2}{c|}{Calculated}&\multicolumn{2}{c}{Difference R—B}\\
1803&\shortstack{Mean time\\in Bremen}&R.A.&Decl.&R.A.&Decl.\\\hline
Feb.\ 21&$17^{\mathrm u}\ 6^{\mathrm m}\ 10^{\mathrm s}$&$272^0\ 58'\ 47''$&$7^0\ 30'\ 40''$&$+2'\ 2''$&$-34''$\\
23&$15\ 24\ 36$&$273\ 31\ 14$&$7\ 45\ 4$&$+2\ 35$&$-57$
\end{tabular}\endgroup
\end{center}

I could not resist making a preliminary correction to the elements on this basis. I shall wait to give better ones until more numerous and more precise observations are available. Your — will see from the comparison with elements V [p.\ 233 of this volume] that the difference is very small:
\SourceLeaf{240}%
\begin{center}
\begingroup\small\setlength{\tabcolsep}{3pt}
\begin{tabular}{l r}
Epoch 1803 on the Seeberg meridian $\ .\ .$&$221^0\ 28'\ 54''0$\\
Annual trop. motion (365 d.) $\ .\ .$&$78\ 0\ 36.9$\\
Daily $\ .\ .\ .\ .\ .\ .\ .\ .\ .\ .\ .$&$769''416$\\
Aphelion 1803, fixed relative to the stars $\ .$&$301^0\ 24'\ 13''$\\[3pt]
Node $\ .\ .\ .\ .\ .\ .\ .\ .\ .\ .\ .$&$172^0\ 28'\ 8''$\\
Eccentricity $\ .\ .\ .\ .\ .\ .\ .\ .\ .$&$0.245619$\\
Logarithm of the semimajor axis&$0.4426160$\\
Inclination of the orbit $\ .\ .\ .\ .\ .\ .\ .\ .$&$34^0\ 38'\ 20''$
\end{tabular}\endgroup
\end{center}

These elements represent the observation as follows:
\begin{center}
\begingroup\small\setlength{\tabcolsep}{4pt}
\begin{tabular}{l|r|r|r|r}
&\multicolumn{2}{c|}{Calculated}&\multicolumn{2}{c}{Difference R—B}\\
1803&R.A.&Decl.&R.A.&Decl.\\\hline
Feb.\ 21&$272^0\ 56'\ 25''0$&$7^0\ 31'\ 29''8$&$-20''0$&$+15''8$\\
23&$273\ 28\ 46.8$&$7\ 45\ 52.3$&$+7.8$&$-7.7$
\end{tabular}\endgroup
\end{center}
It will therefore be possible to find the planet throughout this year as easily as one could wish.

It is to be hoped that $\CeresSymbol$ too will soon be found again. Since it lies low in the southern sky this time, it will probably be found sooner in Italy and France than here.

